% questions.tex —— 1958 年·云南卷 · 题目内容片段（无答案）
% 由 tools/generate_projects.py 依来源总账自 1958/1958_云南_数学真题_试题.md 生成，勿手改（改源 MD 后重跑）


\begin{question}
\gkti{1}解方程 $\left(\log x^{\sqrt5}\right)^2+\log x^5-\dfrac54=0$。

\end{question}
\begin{question}
\gkti{2}求 $(1+x)+(1+x)^2+\cdots+(1+x)^{10}$ 式中 $x^8$ 项的系数（不必求出其它各项的系数）。

\end{question}
\begin{question}
\gkti{3}求圆锥曲线 $2x^2-8xy-4y^2-4y+1=0$ 之焦点及准线。

\end{question}
\begin{question}
\gkti{4}求圆锥曲线 $x^2+y^2=49$ 与 $x^2+y^2-20y+91=0$ 之公切线。

\end{question}
\begin{question}
\gkti{5}证明：设一四面体为一平面所截，若截面是平行四边形，则此截面平行于四面体的两条棱。

\end{question}
\begin{question}
\gkti{6}设三角形 $ABC$ 中，$\angle A$ 的外角平分线与 $BC$ 边之延长线相交于 $D$，与外接圆周交于 $E$，求证：$AB\cdot AC=AE\cdot AD$。

\end{question}
\begin{question}
\gkti{7}设有互相外切的二圆的半径各为 $a$ 与 $b$（$a>b$），作此二圆的两外公切线，设其夹角为 $A$，证明：

\nopagebreak[4]
\[
\sin A=\frac{4(a-b)\sqrt{ab}}{(a+b)^2}
\]

\end{question}
\begin{question}
\gkti{8}在三角形 $ABC$ 中，$r$ 与 $R$ 分别为其内切圆与外接圆的半径，求证：

\nopagebreak[4]
\[
r=4R\sin\frac A2\sin\frac B2\sin\frac C2
\]

\end{question}